\documentclass[reqno, answers]{exam}

\usepackage{amssymb, amsmath, amsthm, stmaryrd}
\usepackage{fullpage, parskip}
\usepackage{graphicx}
\usepackage{enumitem}

% Improved spacing
\usepackage[top=1in, bottom=0.8in, left=1in, right=1in]{geometry}


\title{\normalsize{
  JHED: \underline{\hspace{4cm}} \hspace{1em} Name: \underline{\hspace{9cm}}}
   }
\author{}
\date{}

% Custom spacing for better page utilization
\renewcommand{\baselinestretch}{0.95}

\begin{document}
\maketitle
\vspace{-2cm} % Reduce space after title

\begin{questions}
  \question[8]
Consider a hypothetical machine using \textit{base 10} arithmetic with precision $d = 2$ (two digits in the fractional part). Let $\mathbb{F}$ be the set of floating point numbers representable on this machine.

\begin{parts}
  \part State the machine epsilon $\epsilon_{\text{mach}}$ for this system.
  \begin{solution}
  $\epsilon_{\text{mach}} = \beta^{-d} = 10^{-2} = \boxed{0.01}$.
  \end{solution}
  \vfill
  \part Count the elements of $\mathbb{F}$ in the interval $[40, 4000]$, including both endpoints.
  \begin{solution}
  With base 10 and $d=2$, normalized numbers have the form $a.bc\times 10^n$, with $a\in\{1,\dots,9\}$ and $b,c\in\{0,\dots,9\}$. Within $[40,4000]$ at integer granularity:

  In $[40, 99]$: $40, 41, \ldots, 99$ $\rightarrow$ 60 elements.

  In $[100, 999]$: $100, 101, \ldots, 999$ $\rightarrow$ 900 elements.

  In $[1000, 4000]$: $1000, 1001, \ldots, 4000$ $\rightarrow$ 3001 elements.

  Total: $60 + 900 + 3001 = \boxed{3961}$ elements.
  \end{solution}
  \vfill
  \part Identify the element of $\mathbb{F}$ that is nearest to the real number $e \approx 2.71828$.
  \begin{solution}
    With $d=2$, the nearest is $\boxed{2.72}$.
  \end{solution}
  \vfill
  \part Provide an example of 10 consecutive positive integers that do not belong to $\mathbb{F}$ at some larger magnitude.
  \begin{solution}
   One valid example: $\boxed{100001, 100002, 100003, 100004, 100005, 100006, 100007, 100008, 100009, 100010}$.
  \end{solution}
  \vfill
\end{parts}

  \newpage
  Let $a = 8.4367$, $b = 5.2794$, $c = 1.8649$.

  \question[6]
  Suppose $\mathrm{fl}(x)$ denotes the element of $\mathbb{F}$ \textbf{closest} to the real number $x$ (round half up) under $d=2$ precision.
  Compute the following, and show intermediate steps:

\begin{parts}
  \part $a + (b - c)$
  \begin{solution}
  \textbf{Step 0:} Round inputs to $d=2$: $\mathrm{fl}(a) = 8.44$, $\mathrm{fl}(b) = 5.28$, $\mathrm{fl}(c) = 1.86$.
  
  \textbf{Step 1:} Compute $b - c = 5.28 - 1.86 = 3.42$ (already $d=2$).
  
  \textbf{Step 2:} Compute $a + 3.42 = 8.44 + 3.42 = 11.86$.
  
  \textbf{Step 3:} Round to $d=2$: $\mathrm{fl}(11.86) = 11.86$.
  
  $$\boxed{11.86}$$
  \end{solution}
  \vfill
  \part $(a + b) - c$
  \begin{solution}
  \textbf{Step 0:} Round inputs to $d=2$: $\mathrm{fl}(a) = 8.44$, $\mathrm{fl}(b) = 5.28$, $\mathrm{fl}(c) = 1.86$.
  
  \textbf{Step 1:} Compute $a + b = 8.44 + 5.28 = 13.72$.
  
  \textbf{Step 2:} Round to $d=2$: $\mathrm{fl}(13.72) = 13.72$.
  
  \textbf{Step 3:} Compute $13.72 - c = 13.72 - 1.86 = 11.86$.
  
  \textbf{Step 4:} Round to $d=2$: $\mathrm{fl}(11.86) = 11.86$.
  
  $$\boxed{11.86}$$
  \end{solution}
  \vfill
\end{parts}

  \question[6]
  Now suppose $\mathrm{fl}(x)$ is obtained by truncation (chopping) to $d=2$ precision, i.e., remove all digits beyond two fractional digits.
  Compute the following, and show intermediate steps:

\begin{parts}
  \part $a + (b - c)$
  \begin{solution}
  \textbf{Step 0:} Truncate inputs to $d=2$: $\mathrm{fl}(a) = 8.43$, $\mathrm{fl}(b) = 5.27$, $\mathrm{fl}(c) = 1.86$.
  
  \textbf{Step 1:} Compute $b - c = 5.27 - 1.86 = 3.41$ (already $d=2$).
  
  \textbf{Step 2:} Compute $a + 3.41 = 8.43 + 3.41 = 11.84$.
  
  \textbf{Step 3:} Truncate to $d=2$: $\mathrm{fl}(11.84) = 11.84$.
  
  $$\boxed{11.84}$$
  \end{solution}
  \vfill
  \part $(a + b) - c$
  \begin{solution}
  \textbf{Step 0:} Truncate inputs to $d=2$: $\mathrm{fl}(a) = 8.43$, $\mathrm{fl}(b) = 5.27$, $\mathrm{fl}(c) = 1.86$.
  
  \textbf{Step 1:} Compute $a + b = 8.43 + 5.27 = 13.70$.
  
  \textbf{Step 2:} Truncate to $d=2$: $\mathrm{fl}(13.70) = 13.70$.
  
  \textbf{Step 3:} Compute $13.70 - c = 13.70 - 1.86 = 11.84$ (already $d=2$).
  
  $$\boxed{11.84}$$
  \end{solution}
  \vfill
\end{parts}

\end{questions}
\end{document}